% 中文稿件独立入口。请使用 XeLaTeX 编译：
% latexmk -xelatex -interaction=nonstopmode -halt-on-error paper_zh.tex
\documentclass[twocolumn]{hust}

\usepackage[toc,page,header]{appendix}
\usepackage{minitoc}
\usepackage{amsmath,amssymb}
\usepackage{bm}
\usepackage{wrapfig}
\usepackage{enumitem}
\usepackage{xeCJK}

% 使用 TeX Live 自带的 Fandol 字体，避免依赖本机安装的中文字体。
\setCJKmainfont{FandolSong-Regular.otf}[
  BoldFont=FandolSong-Bold.otf
]
\setCJKsansfont{FandolHei-Regular.otf}[
  BoldFont=FandolHei-Bold.otf
]
\setCJKmonofont{FandolFang-Regular.otf}
\renewcommand{\hustcorrespondencelabel}{通讯作者}
\renewcommand{\hustreportlabel}{技术报告}

\addto{\captionsenglish}{%
  \renewcommand{\contentsname}{目录}%
  \renewcommand{\figurename}{图}%
  \renewcommand{\tablename}{表}%
  \renewcommand{\refname}{参考文献}%
}
\renewcommand{\beginabstract}{\abstractfont{\sffamily\bfseries\hustblue{摘要}\par}}
\renewcommand{\beginappendix}{\appendix{\titlefont\sffamily\hustblue{附录}\par}}
\algrenewcommand{\algorithmicrequire}{\sffamily\bfseries\color{hustblue}输入：}
\algrenewcommand{\algorithmicensure}{\sffamily\bfseries\color{hustblue}输出：}

\newcommand{\templatefigure}[2][0.95\linewidth]{%
  \fbox{%
    \begin{minipage}[c][0.28\linewidth][c]{#1}
      \centering\sffamily\small
      #2\\[0.6em]
      请替换为正式图片或表格。
    \end{minipage}%
  }%
}

\title{论文标题：简洁概括核心贡献的副标题}

\author[1]{第一作者\textsuperscript{*}}
\author[1,2]{第二作者\textsuperscript{*}}
\author[2]{第三作者}

\affiliation[1]{华中科技大学 XXX 学院}
\affiliation[2]{华中科技大学 YYY 学院}

\email{author1@hust.edu.cn}
\email{author2@hust.edu.cn}
\email{author3@hust.edu.cn}
\contribution{\textsuperscript{*} 共同第一作者}
\correspondence{author1@hust.edu.cn}

\abstract{
请在此撰写中文摘要，简要说明研究问题、方法、主要贡献和实验结论。建议用后续句子补充核心技术思路、评估设置以及最重要的结果或发现。
}

\checkdata[资源链接]{%
  \href{https://choucisan.github.io}{\raisebox{-1ex}{\includegraphics[height=4.2ex]{assets/internet-icon.pdf}}~项目主页}\hspace{2em}
  \href{https://github.com/choucisan}{\raisebox{-1ex}{\includegraphics[height=4.2ex]{assets/github-logo.pdf}}~代码}\hspace{2em}
  \href{https://huggingface.co/choucsan}{\raisebox{-1ex}{\includegraphics[height=4.2ex]{assets/hf-logo.pdf}}~模型}
}

\begin{document}
\ifhusttwocolumn
  \renewcommand{\hustfrontmatter}{%
    \par\vskip 2mm
    \centering
    \templatefigure[0.92\textwidth]{论文概览占位图}
    \captionof{figure}{\textbf{论文概览。} 请替换为论文的主要示意图、流程图或关键结果图。}
    \label{fig:teaser}%
  }
\fi
\maketitle
\vspace{-1.1cm}

\ifhusttwocolumn
  \onecolumn
  \vspace*{-1.1cm}
  \begin{center}
    \tableofcontents
  \end{center}
  \clearpage
  \twocolumn
\else
  \begin{figure}[!htbp]
    \centering
    \templatefigure[0.98\linewidth]{论文概览占位图}
    \caption{\textbf{论文概览。} 请替换为论文的主要示意图、流程图或关键结果图。}
    \label{fig:teaser}
  \end{figure}
  \tableofcontents
  \newpage
\fi

% 中文正文采用 1.5 倍行距；标题页和目录保持模板原有行距。
\onehalfspacing

\section{引言}

介绍研究问题及其意义。本节开头应说明应用场景、现有工作的局限，以及开展本研究的实际或科学动机。

概括本文要解决的核心问题。问题可以涉及质量、效率、鲁棒性、可扩展性、可解释性、数据可用性或部署约束。

简要介绍本文方法，为读者建立整体认识。适当时可引用论文概览图，例如图~\ref{fig:teaser}。

本文的主要贡献如下：
\begin{itemize}
  \item \hustred{贡献一。} 用一至两句话描述第一项技术或实证贡献。
  \item \hustred{贡献二。} 说明第二项创新及其价值。
  \item \hustred{贡献三。} 说明验证、数据集、基准、应用或分析方面的贡献。
\end{itemize}

最后简要总结支撑上述主张的证据，包括评估设置、关键指标和总体结论；尚未核实的数值不应提前写入。

\newpage

\section{方法}

介绍所提出的方法，使读者能够理解并复现核心思路。先概述整体流程，再按处理顺序说明各组成部分。

\subsection{问题定义}

定义输入、输出、假设和符号。例如，令 $x$ 表示输入，$y$ 表示目标输出，$f_\theta$ 表示待优化的模型或算法。说明任务属于监督、自监督、生成式、交互式还是分析型设置。

\subsection{系统概览}

\begin{figure}[t]
  \centering
  \templatefigure[0.95\linewidth]{方法概览占位图}
  \caption{\textbf{方法概览。} 请替换为展示输入、主要模块和输出的流程图。}
  \label{fig:method_overview}
\end{figure}

结合图~\ref{fig:method_overview} 说明系统的主要组成部分。本小节侧重整体概念，架构、损失函数和实现细节可放在后续小节。

\subsection{核心模块}

说明核心技术模块解决的问题、接收的信息以及产生的表示或预测。若使用神经网络，介绍骨干网络、条件信号和输出头；若使用其他方法，说明算法阶段和决策规则。

\subsection{优化目标}

给出训练或优化目标，并在公式之后立即定义各项：
\begin{equation}
  \mathcal{L} = \lambda_1 \mathcal{L}_{\mathrm{main}} + \lambda_2 \mathcal{L}_{\mathrm{reg}}.
\end{equation}
其中，$\mathcal{L}_{\mathrm{main}}$ 是主要任务损失，$\mathcal{L}_{\mathrm{reg}}$ 是可选的正则项。

\subsection{实现细节}

总结复现所需的实现选择，包括模型规模、优化器、学习率计划、批大小、硬件和预处理方式。较长的细节可移至附录。

\section{数据}
\label{sec:data}

\begin{figure*}[t]
  \centering
  \templatefigure[0.92\linewidth]{数据流程占位图}
  \caption{\textbf{数据流程。} 请替换为数据集构建、预处理或标注流程图。}
  \label{fig:data_pipeline}
\end{figure*}

介绍数据来源、收集流程、筛选标准、标注过程和数据划分，使读者了解数据范围及其局限，同时避免泄露未公开或隐私信息。

\textbf{数据收集。} 说明数据来源及纳入标准；适用时报告样本数、领域、类别、标注人员或采集设备。

\textbf{预处理。} 描述清洗、归一化、去重、质量控制和格式转换过程。若展示了处理流程，可引用图~\ref{fig:data_pipeline}。

\textbf{划分与使用。} 说明训练集、验证集和测试集的构造方式，以及防止数据泄漏的措施。

\textbf{伦理与许可。} 说明适用的知情同意、隐私、许可和安全考量。对于公开数据集，请引用其原始来源，例如占位引用~\cite{sample_reference}。

\section{训练}

\subsection{训练设置}

描述主要模型或算法的训练方案，包括优化器、学习率计划、批大小、训练步数或轮数、硬件、精度和检查点选择策略。

\subsection{训练阶段}

如果方法包含多个阶段，请按时间顺序描述：
\begin{itemize}
  \item \textbf{阶段一。} 说明第一阶段的目标和使用的数据。
  \item \textbf{阶段二。} 说明第二阶段在目标函数、数据、分辨率或监督信号上的变化。
  \item \textbf{最终调优。} 说明最后的适配、校准或模型选择步骤。
\end{itemize}

\subsection{开发过程中的消融}

简要说明开发过程中测试过的设计选择。详细消融结果可放在第~\ref{sec:evaluation} 节或附录中。

\section{推理与部署}

\subsection{推理流程}

说明训练完成后模型或方法在测试阶段的使用方式，包括输入、预处理、解码或后处理步骤以及预期输出。如果推理过程与训练过程不同，请明确说明。

\begin{algorithm}[t]
  \caption{\textbf{推理流程示例。} 请替换为实际算法步骤。}
  \label{alg:inference}
  \begin{algorithmic}[1]
    \Require 输入样本 $x$，训练完成的模型 $f_{\theta}$
    \Ensure 预测结果 $\hat{y}$
    \State $z \gets \textsc{Preprocess}(x)$
    \State $s \gets f_{\theta}(z)$
    \If{$\textsc{Confidence}(s) < \tau$}
      \State $s \gets \textsc{Refine}(z, s)$
    \EndIf
    \State $\hat{y} \gets \textsc{Decode}(s)$
    \State \Return $\hat{y}$
  \end{algorithmic}
\end{algorithm}

\subsection{效率}

在相关时报告运行时间、内存占用、模型大小或计算复杂度。比较多个方案时可使用表格。最终基准数据核实前，避免写入未经确认的结论。

\subsection{实践注意事项}

说明实际使用中的限制，例如输入分辨率、支持格式、批处理行为、失败情况或与下游工具的兼容性。

\section{实验}
\label{sec:evaluation}

本节用于验证论文的主要主张。首先说明评估任务、数据集、指标和基线方法。

\begin{figure*}[!h]
  \centering
  \templatefigure[0.9\linewidth]{定性对比占位图}
  \caption{\textbf{定性对比。} 请替换为有代表性的输出、样例或可视化对比。}
  \label{fig:qualitative_results}
\end{figure*}

\subsection{实验设置}

描述基准协议、基线、指标和统计检验流程。明确哪些结果来自本文方法，哪些是复现自已有工作。

\subsection{主要结果}

报告主要定量结果。数值对比可使用表格，趋势或定性行为可使用图示，并结合引言中的研究主张讨论结果。

\begin{table}[t]
  \centering
  \caption{\textbf{结果示例。} 请将占位数值替换为主要对比结果。}
  \label{tab:results}
  \small
  \setlength{\tabcolsep}{5pt}
  \begin{tabular}{lccc}
    \toprule
    {\sffamily\bfseries\color{hustblue}方法} &
    {\sffamily\bfseries\color{hustblue}准确率} &
    {\sffamily\bfseries\color{hustblue}F1} &
    {\sffamily\bfseries\color{hustblue}提升} \\
    \midrule
    基线 A & 78.4 & 75.9 & -- \\
    基线 B & 80.1 & 77.6 & +1.7 \\
    本文方法 & \hustred{83.6} & \hustred{81.4} & \hustred{+3.5} \\
    \bottomrule
  \end{tabular}
\end{table}

\subsection{消融实验}

分析关键设计选择的作用。每组消融实验都应回答一个明确问题。

\subsection{局限性}

说明已知局限、失败案例及不适用的条件，以便读者正确理解结果。

\section{应用}

介绍所提出方法在实际工作流程中的使用方式。本节为可选内容，适用于具有基准分数以外下游价值的工作。

\subsection{应用场景一}

\begin{figure}[t]
  \centering
  \templatefigure[0.95\linewidth]{应用示例占位图}
  \caption{\textbf{应用示例。} 请替换为展示方法如何用于下游场景的示例。}
  \label{fig:application_example}
\end{figure}

说明任务、用户流程、集成位置和预期收益。若图中展示了应用流程，可引用图~\ref{fig:application_example}。

\subsection{应用场景二}

描述第二种用例或部署环境，并尽可能联系第~\ref{sec:evaluation} 节报告的可测量指标。

\section{结论}

用一段简洁文字总结本文问题、方法和主要发现。避免在结论中引入新的技术细节。

随后讨论工作的更广泛意义、局限性和未来方向。可在此说明计划中的扩展、尚待解决的问题或负责任部署方面的考量。

\newpage


\clearpage
\bibliographystyle{IEEEtran}
\bibliography{main_zh}

\clearpage
\beginappendix
\section{补充细节}
\label{sec:appendix_details}

附录用于放置支撑论文结论、但会打断正文叙述的材料，例如扩展实现细节、补充实验、证明、额外定性示例、标注说明和数据集文档。

\subsection{可复现性清单}

\begin{itemize}
  \item 列出软件依赖及其版本。
  \item 说明训练和评估使用的硬件。
  \item 记录随机种子和评估协议。
  \item 发布代码、模型和数据时，提供获取链接或使用说明。
\end{itemize}

\subsection{补充结果}

在此添加补充表格、图表和定性示例。


\end{document}
